\chapter{Momentum, impulse and collisions}\label{ch:momentum}

Force is what changes momentum: $\mathbf F=d\mathbf p/dt$ is Newton's
second law in its original form, and the reason $\mathbf p=m\mathbf v$ is
the right bookkeeping unit for a chapter about things that hit each other.
Momentum is the first \emph{conserved} quantity in the book --- the thing
that survives collisions unchanged --- and the discipline of \emph{when}
it survives (the system-boundary rule) is the real lesson.

\section{Momentum and impulse}\label{sec:mom-p}

$\mathbf p=m\mathbf v$ is a vector. Impulse\index{impulse} is momentum's
denomination of payment:
\begin{equation}
  \mathbf J=\int\mathbf F\,dt=\Delta\mathbf p.\label{eq:impulse}
\end{equation}
A cricket ball stopped in a glove and a foam wall deliver the same $\Delta p$
--- impulse spreads the payment over different times (figure~\ref{fig:impulse}),
which is why crumple zones, airbags, and bent knees all buy
\emph{time}.

\begin{figure}[htbp]\centering
\begin{tikzpicture}
\begin{axis}[width=8.5cm,height=5cm,xlabel={$t$ (ms)},ylabel={$F$ (kN)},
xmin=0,xmax=14,ymin=0,ymax=9,xtick={0,2,4,6,8,10,12,14},font=\small]
\addplot[very thick,accent,domain=0:10] {8*sin(deg(pi*x/10))};
\addplot[fill=accent!15,draw=none,domain=0:10] {8*sin(deg(pi*x/10))} \closedcycle;
\node[font=\scriptsize] at (axis cs:5,4) {$J=\displaystyle\int F\,dt$};
\draw[->,thick,amber] (axis cs:11,6) -- (axis cs:13,6) node[right,font=\scriptsize]{$\Delta\mathbf p$};
\end{axis}
\end{tikzpicture}

\caption{Impulse is area under the $F(t)$ curve. Notice the same shaded
area can be a tall narrow spike or a long gentle hill --- same $\Delta p$,
very different peak force.}
\label{fig:impulse}
\end{figure}

\begin{example}[Impulse from a real $F(t)$ curve]\label{ex:impulse}
A $0.15$\,kg baseball hits a bat with the half-sine force profile
$F(t)=8000\sin(\pi t/10\,\mathrm{ms})$\,N. Find the impulse and the exit
speed if the ball arrived at $40\,\mathrm{m/s}$.
\par\smallskip
\textbf{Strategy.} Area under the curve: the half-sine's average height is
$2/\pi$ of peak over its $10$\,ms width.
\begin{solution}
$J=\int_0^{10}F\,dt=\frac{2}{\pi}\times8000\times0.010=50.9$\,N\,s.
Reversing the ball needs $\Delta p=2mv=2\times0.15\times40=12$\,N\,s ---
the bat over-supplies: exit speed satisfies
$m(v_{\rm out}+40)=50.9$, giving $v_{\rm out}=299$\,m/s. (Real bats don't:
much of $J$ deforms the ball, so effective impulse is lower --- but the
method is the point.) \checkmark
\end{solution}
\end{example}

\section{Systems: the boundary is a choice}\label{sec:mom-sys}

Internal forces come in third-law pairs and cancel inside the sum; only
\emph{external} forces move the total momentum. Drawing the dashed line of
figure~\ref{fig:boundary} is a \emph{decision}: put both blocks inside and
the contact force vanishes from the bookkeeping; put it outside and the
same force is external and changes $\mathbf p$.

\begin{figure}[htbp]\centering
\begin{tikzpicture}[scale=1.0,force/.style={->,thick,accent}]
\filldraw[fill=gray!20] (0,0) rectangle (1,0.7);
\filldraw[fill=gray!20] (1.6,0) rectangle (2.6,0.7);
\draw[thick] (-0.4,-0.12) -- (3.0,-0.12);
\node[font=\scriptsize] at (0.5,0.35) {$m_1$};
\node[font=\scriptsize] at (2.1,0.35) {$m_2$};
\draw[force] (1,0.35) -- (1.6,0.35) node[midway,above,font=\scriptsize]{$\mathbf F_{12}$};
\draw[force] (1.6,0.35) -- (1,0.35) node[midway,below,font=\scriptsize]{$\mathbf F_{21}$};
\draw[dashed,thick,amber] (-0.3,-0.3) rectangle (2.9,1.0);
\node[font=\scriptsize,amber] at (1.3,1.15) {system: both blocks --- $F_{12},F_{21}$ internal};
\begin{scope}[xshift=4.2cm]
\filldraw[fill=gray!20] (0,0) rectangle (1,0.7);
\filldraw[fill=gray!20] (1.6,0) rectangle (2.6,0.7);
\draw[thick] (-0.4,-0.12) -- (3.0,-0.12);
\draw[force] (1,0.35) -- (1.6,0.35) node[midway,above,font=\scriptsize]{$\mathbf F_{12}$};
\draw[dashed,thick,amber] (-0.3,-0.3) rectangle (1.3,1.0);
\node[font=\scriptsize,amber] at (1.3,1.15) {system: $m_1$ only --- $F_{12}$ external};
\end{scope}
\end{tikzpicture}

\caption{The same collision under two system choices. Internal forces
cancel by the third law; external ones change the system's $\mathbf p$.}
\label{fig:boundary}
\end{figure}

\section{Conservation of momentum: the applicability rule}\label{sec:mom-cons}

\begin{principle}[Momentum conservation condition]
$\sum\mathbf p$ is conserved for a system exactly when the \emph{net
external impulse} during the event is zero. Gravity during a $10$\,ms
collision contributes $\sim10^{-3}$\,N\,s --- negligible. Gravity during a
$10$\,s flight is not.
\end{principle}

This is why ``momentum is conserved in collisions'' has the silent rider
``while the collision lasts''. The ballistic pendulum's swing, the
projectile's arc --- neither conserves $\mathbf p$; the \emph{impact} does.

\section{Centre of mass: the system's average representative}\label{sec:mom-cm}

\[
  \mathbf r_{\rm cm}=\frac{1}{M}\sum_i m_i\mathbf r_i,\qquad
  M\mathbf A_{\rm cm}=\mathbf F_{\rm ext}.
\]
The CM\index{centre of mass} obeys single-particle mechanics under the
external force alone --- internal explosions, splits, and collisions can't
move it (figure~\ref{fig:cm}). A wrench thrown across the room tumbles in
a mess; its CM traces the clean parabola of chapter~1.

\begin{figure}[htbp]\centering
\begin{tikzpicture}[scale=1.0]
\draw[thick] (-0.5,0) -- (7.5,0);
\foreach \x in {0,1,...,7} \draw (\x,0.05) -- (\x,-0.05) node[below,font=\scriptsize]{\x};
\filldraw[fill=accent!30] (1.8,0) rectangle (2.4,0.5);
\node[font=\scriptsize] at (2.1,0.7) {$m_1=1$ kg};
\filldraw[fill=accent!30] (5.8,0) rectangle (6.4,0.5);
\node[font=\scriptsize] at (6.1,0.7) {$m_2=3$ kg};
\draw[->,very thick,amber] (5.2,0.25) -- (5.2,1.1) node[above,font=\scriptsize]{CM $=5.2$ m};
\node[font=\scriptsize] at (3.6,-0.6) {$x_{\rm cm}=\frac{m_1x_1+m_2x_2}{m_1+m_2}$};
\end{tikzpicture}

\caption{CM of two masses: $x_{\rm cm}$ weights each position by mass.
Notice it sits nearer the heavy block --- CM is the mass-weighted average,
not the midpoint.}
\label{fig:cm}
\end{figure}

\begin{example}[Person on a movable cart]\label{ex:cart}
A $60$\,kg person stands on a $20$\,kg cart (wheels frictionless) and walks
$3$\,m forward along the cart. How far does the cart move under them?
\par\smallskip
\textbf{Strategy.} No external horizontal force $\Rightarrow$ the CM of
person$+$cart doesn't move. The person moves $3$\,m relative to the
\emph{cart}, not the ground --- the classic mistake.
\begin{solution}
Let the cart move $-x$ (backward) while the person moves $3-x$ forward on
the ground. CM fixed:
$60(3-x)+20(-x)=0 \Rightarrow 180=80x \Rightarrow x=2.25$\,m.
The cart slides back $2.25$\,m; the person gains only $0.75$\,m of ground.
This is why stepping off a light boat onto a dock is hazardous. \checkmark
\end{solution}
\end{example}

\begin{exercise}[easy]
Redo example~\ref{ex:cart} with a $90$\,kg person and a $10$\,kg cart ---
does the cart move more or less? Compute it. (AP-C)
\end{exercise}

\section{1-D collisions: the two canonical cases}\label{sec:mom-1d}

\textbf{Perfectly inelastic} (they stick): momentum alone fixes the
answer --- the shared speed is total momentum over total mass,
$v'=(m_1v_1+m_2v_2)/(m_1+m_2)$ --- and $K$ drops by the fraction that
went into deformation. \textbf{Elastic}: momentum plus energy conservation
(or equivalently, relative velocity reverses). The full elastic solution
needs $K$, so \S4.7 of the next chapter closes it; this chapter keeps the
momentum half. Figure~\ref{fig:collseq} shows the three phases --- and that
the force between the blocks is internal either way.

\begin{figure}[htbp]\centering
\begin{tikzpicture}[scale=1.0,force/.style={->,thick,accent}]
% before
\begin{scope}
\filldraw[fill=gray!20] (0,0) rectangle (0.9,0.6);
\filldraw[fill=gray!20] (2.5,0) rectangle (3.4,0.6);
\draw[force] (0.9,0.3) -- (1.7,0.3) node[right,font=\scriptsize]{$v$};
\node[font=\scriptsize] at (0.45,0.85) {$m$}; \node[font=\scriptsize] at (2.95,0.85) {$m$};
\draw[thick] (-0.3,-0.1) -- (3.9,-0.1);
\node[font=\scriptsize] at (1.7,-0.45) {before};
\end{scope}
% during
\begin{scope}[yshift=-1.7cm]
\filldraw[fill=gray!20] (1.2,0) rectangle (2.1,0.6);
\filldraw[fill=gray!20] (2.1,0) rectangle (3.0,0.6);
\draw[force] (3.0,0.3) -- (3.8,0.3) node[right,font=\scriptsize]{$v'$};
\draw[thick] (-0.3,-0.1) -- (4.6,-0.1);
\node[font=\scriptsize] at (1.7,-0.45) {during: $F$ acts equally, oppositely};
\end{scope}
% after elastic
\begin{scope}[yshift=-3.4cm]
\filldraw[fill=gray!20] (1.2,0) rectangle (2.1,0.6);
\filldraw[fill=gray!20] (3.4,0) rectangle (4.3,0.6);
\draw[force] (4.3,0.3) -- (5.1,0.3) node[right,font=\scriptsize]{$v$};
\draw[thick] (-0.3,-0.1) -- (5.6,-0.1);
\node[font=\scriptsize] at (1.7,-0.45) {after ($e=1$): velocities exchanged};
\end{scope}
\end{tikzpicture}

\caption{Collision sequence. The interaction force is internal to the
system whatever the elasticity --- $p$ survives, $K$ doesn't always.}
\label{fig:collseq}
\end{figure}

\begin{checkpoint}[Quick K audit]
Equal masses, one moving: perfectly inelastic collision loses exactly half
the kinetic energy. Run the numbers yourself before the code below does.
\end{checkpoint}

\section{2-D collisions and explosions}\label{sec:mom-2d}

In two dimensions, momentum conservation is two scalar equations --- $x$
and $y$ separately (figure~\ref{fig:glance}). Explosions are collisions run
backward: a stationary object splitting into fragments has
$\sum\mathbf p=0$ after as before, so the pieces fly apart with their
momenta summing to zero --- heavier fragments go slower.

\begin{figure}[htbp]\centering
\begin{tikzpicture}[scale=1.0,force/.style={->,thick,accent}]
% momentum conservation in 2D: p_in = p1' + p2'
\draw[force,very thick] (0,0) -- (3.2,0) node[midway,below,font=\scriptsize]{$\mathbf p_{\rm in}$};
\draw[force,thick,amber] (0,0.08) -- (2.3,1.05) node[right,font=\scriptsize]{$\mathbf p_1'$};
\draw[force,thick,amber] (2.3,1.05) -- (3.2,0) node[midway,right,font=\scriptsize]{$\mathbf p_2'$};
\draw[dashed] (2.3,1.05) -- (2.3,0) node[below,font=\scriptsize]{};
\node[font=\scriptsize] at (0.9,-0.55) {$x$: conserved};
\node[font=\scriptsize] at (2.7,0.65) {$y$: sums to zero};
\end{tikzpicture}

\caption{A 2-D collision as a vector triangle:
$\mathbf p_{\rm in}=\mathbf p_1'+\mathbf p_2'$. Notice the $y$-components
must cancel --- perpendicular momentum can't appear from nothing.}
\label{fig:glance}
\end{figure}

\begin{example}[Glancing collision]\label{ex:glance}
A puck at $6$\,m/s strikes an identical stationary puck and leaves at
$4$\,m/s, $30^{\circ}$ above the original direction. Find the second puck's
velocity.
\par\smallskip
\textbf{Strategy.} Equal masses: velocity conservation is vector
conservation. Components: $x$ must return $6$, $y$ must net $0$.
\begin{solution}
$x$: $4\cos30^{\circ}+v_{2x}=6 \Rightarrow v_{2x}=6-3.464=2.54$\,m/s.
$y$: $4\sin30^{\circ}+v_{2y}=0 \Rightarrow v_{2y}=-2$\,m/s.
$|\mathbf v_2|=\sqrt{2.54^{2}+4}=3.2$\,m/s at $-38.2^{\circ}$. The angle
sum $30+38.2\approx68^{\circ}$ --- not the elastic equal-mass
$90^{\circ}$, so this collision lost energy. \checkmark
\end{solution}
\end{example}

\begin{exercise}[medium]
A $4$\,kg firework at rest explodes into three pieces: $1$\,kg north at
$12$\,m/s, $1.5$\,kg east at $8$\,m/s. Find the third piece's velocity.
(F=ma)
\end{exercise}

\section{Continuous mass transfer: the rocket equation}\label{sec:mom-rocket}

When mass leaves continuously, momentum bookkeeping over one ejection
(figure~\ref{fig:rocket}) gives
\[
  0=(m-\Delta m)\,\Delta v-u\,\Delta m
  \;\xrightarrow{\ \Delta m\to0\ }\;
  m\,dv=-u\,dm
  \;\Rightarrow\;
  \Delta v=u\ln\frac{m_0}{m_f},
\]
the \emph{Tsiolkovsky rocket equation}\index{rocket equation} --- $u$ the
exhaust speed relative to the rocket. The logarithm is why rockets are
mostly fuel tank: each equal $\Delta v$ costs a multiplicative mass ratio.

\begin{figure}[htbp]\centering
\begin{tikzpicture}[scale=1.0,force/.style={->,thick,accent}]
% rocket
\draw[thick] (0,0) -- (0,2.2) -- (0.4,2.8) -- (0.8,2.2) -- (0.8,0) -- cycle;
\draw[thick] (0,0) -- (-0.4,-0.8); \draw[thick] (0.8,0) -- (1.2,-0.8);
\draw[fill=amber!40] (0.12,-0.1) -- (0.4,-1.3) -- (0.68,-0.1) -- cycle;
\node[font=\scriptsize] at (0.4,1.4) {$m-\Delta m$};
\draw[force] (0.4,-1.3) -- (0.4,-2.2) node[below,font=\scriptsize]{$\Delta m$ at $u$};
\draw[force,amber] (1.2,1.4) -- (1.2,2.6) node[right,font=\scriptsize]{$\Delta v$};
\node[font=\scriptsize,align=left] at (3.6,0.4) {momentum books:\\[2pt]
$0=(m-\Delta m)\Delta v - u\,\Delta m$\\[2pt]
$\Rightarrow\ \Delta v = u\ln\!\frac{m_0}{m_f}$};
\end{tikzpicture}

\caption{One ejection event: throw $\Delta m$ backward at speed $u$ and
conservation of $\mathbf p$ pushes the rest forward.}
\label{fig:rocket}
\end{figure}

\begin{example}[Staging estimate]\label{ex:rocket}
A rocket exhausts at $u=3$\,km/s. (a) Single stage to orbit
($\Delta v\approx9$\,km/s): mass ratio? (b) Two equal stages?
\par\smallskip
\textbf{Strategy.} $m_0/m_f=e^{\Delta v/u}$ each stage; stages multiply.
\begin{solution}
(a) $e^{9/3}=e^{3}\approx20$: the rocket is $95\%$ propellant.
(b) Each stage needs $\Delta v=4.5$\,km/s:
$e^{1.5}\approx4.48$ per stage --- payload structure is dramatically
better. This arithmetic, not optimism, is why real rockets stage. \checkmark
\end{solution}
\end{example}

\begin{example}[Water-jet thrust]\label{ex:jet}
A fire hose delivers $20$\,L/s of water ($\rho=1$\,kg/L) at
$15$\,m/s. What force does the jet exert on a wall it splashes against?
\par\smallskip
\textbf{Strategy.} Momentum flux: force equals mass flow rate times the
speed destroyed, $F=\dot m v$.
\begin{solution}
$F=(20\,\mathrm{kg/s})(15\,\mathrm{m/s})=300\,\mathrm{N}$ --- a
refrigerator's weight, and the reason hoses need two operators. The same
$\dot m v$ formula gives a rocket's thrust before the logarithm does the
integration. \checkmark
\end{solution}
\end{example}

\begin{lstlisting}[language=Python,caption={\texttt{code/ch03\_collision\_chain.py} --- event-stepped collisions: $p$ survives, $K$ audits the loss.},label={lst:coll}]
def collide(m1, v1, m2, v2, e):
    v1p = (m1*v1 + m2*v2 - m2*e*(v1-v2)) / (m1+m2)
    v2p = (m1*v1 + m2*v2 + m1*e*(v1-v2)) / (m1+m2)
    return v1p, v2p
\end{lstlisting}

Output of a real run:
\begin{lstlisting}[numbers=none]
initial:            p= 4.000 kg m/s   K= 8.000 J
after e=0 (1kg+2kg): p= 4.000 kg m/s   K= 2.667 J  (K lost 5.333 J = 66.7%)
after e=1 (3kg+3kg): p= 4.000 kg m/s   K= 2.667 J  v1'=+0.000 v2'=+1.333
e=0.0: va'=+2.000 vb'=+2.000  K=8.000 J (lost 50.0%)
e=0.5: va'=+1.000 vb'=+3.000  K=10.000 J (lost 37.5%)
e=1.0: va'=+0.000 vb'=+4.000  K=16.000 J (lost 0.0%)
\end{lstlisting}

Read it in order: momentum is conserved to machine precision in every line;
kinetic energy is not (the $66.7\%$ loss of the perfectly inelastic hit is
the chapter's promise made numeric). The elastic equal-mass exchange
($e=1$: velocities swap) and the intermediate loss at $e=0.5$ complete the
restitution family.

\section{In engineering practice: the crash budget}\label{sec:mom-eng}

Vehicle crash engineering is impulse management: the occupant's $\Delta p$
is fixed by the speed, so design maximises collision \emph{time} ---
crumple zones, airbags, and collapsible steering columns all stretch
$\Delta t$ to shrink $F$. The same bookkeeping runs in reverse for
spacecraft: a Falcon 9's landing burn spends $\dot m v$ of thrust exactly
as example~\ref{ex:jet} computes, and mass-ratio arithmetic of
example~\ref{ex:rocket} is why propellant dominates every launch manifest.

\section{Common mistakes}\label{sec:mom-mistakes}

\begin{tripwire}[Conserving $p$ when an external impulse exists]
Momentum conservation needs \emph{zero external impulse for the duration} ---
gravity during a long swing, friction during a long skid, and the dock
during a boat step-off all break it.
\end{tripwire}
\begin{tripwire}[Wrong system boundary]
If the fragment you need sits outside your dashed line, conservation
``fails'' --- draw the boundary to include every piece of the event.
\end{tripwire}
\begin{tripwire}[Sign errors in 1-D]
$\mathbf p$ is signed. A rebound at $-v$ contributes
$\Delta p=2mv$, not zero.
\end{tripwire}
\begin{tripwire}[Believing $K$ survives inelastic collisions]
Sticking is maximally lossy --- the equal-mass case loses half. $K$ audits,
$\mathbf p$ conserves.
\end{tripwire}
\begin{tripwire}[Thinking the CM pauses for explosions]
The CM keeps its velocity through the breakup --- the fragments' weighted
sum never moved.
\end{tripwire}

\section{Summary}\label{sec:mom-summary}

\begin{table}[ht]\centering
\begin{tabularx}{\linewidth}{@{}lY@{}}
\toprule
\tablehead\thead{Idea} & \thead{One-line version}\\
\midrule
impulse & $\mathbf J=\int F\,dt=\Delta\mathbf p$ --- area under $F(t)$ \\
boundary rule & internal forces cancel; $\mathbf p$ conserved iff external impulse $=0$ \\
centre of mass & $M\mathbf A_{\rm cm}=\mathbf F_{\rm ext}$; internal events can't move it \\
collisions 1-D & inelastic: $p$ only; elastic: $p+K$ (\S4.7) \\
collisions 2-D & conserve $x$ and $y$ separately \\
rocket & $\Delta v=u\ln(m_0/m_f)$ --- staging is arithmetic \\
thrust & $F=\dot m v$ for continuous flux \\
\bottomrule
\end{tabularx}
\caption{Momentum cheat sheet.}
\label{tab:momsummary}
\end{table}

\begin{wins}[Definition of done]
You can now: compute impulse from a force--time curve; choose a system
boundary and state when $\mathbf p$ is conserved; move a cart under a
walking person without moving the CM; solve inelastic and 2-D collisions;
derive the rocket equation; and read a restitution sweep's $K$ losses.
\end{wins}

\chapternotes{%
The momentum-first presentation follows MIT 8.01SC weeks 3 and 5--6
\cite{ocw801} and AP Physics C unit 4 \cite{apced}. OpenStax College
Physics ch.~8--9 \cite{openstax1} covers collisions and CM at this level;
Morin ch.~5 \cite{morin} is the source of the rocket staging and
variable-mass exercises. The restitution-sweep numbers are the listing's
real output; the baseball numbers are representative.}

\begin{rulebox}[{Watch instead}]
\begin{enumerate}
\item \emph{Momentum, Conservation of Momentum, Center of Mass} --- Walter
Lewin, MIT 8.01x Lecture 15, 52:17, \url{https://youtu.be/aIhScO3_I50} ---
\S\ref{sec:mom-p}--\S\ref{sec:mom-cm} at full length (accessed 2026-10-11).
\item \emph{Elastic and inelastic collisions} --- Khan Academy Physics,
14:24, \url{https://youtu.be/8OB8eIPgEkQ} ---
\S\ref{sec:mom-1d}--\S\ref{sec:mom-2d} worked quickly (accessed 2026-10-11).
\end{enumerate}
\end{rulebox}

\section*{Further reading}
\addcontentsline{toc}{section}{Further reading}

\begin{itemize}
\item MIT OpenCourseWare, \emph{8.01SC Classical Mechanics}, weeks 3, 5--6
--- momentum and collisions with recitation solutions. At
\url{https://ocw.mit.edu}, course ``8.01SC'' (verified 2026-10-11).
\item OpenStax, \emph{College Physics 2e}, ch.~8 ``Linear Momentum and
Collisions'' --- free text, the same laws with more worked explosions. At
\url{https://openstax.org} (verified 2026-10-11).
\item College Board, \emph{AP Physics C: Mechanics CED}, Unit 4 --- the
exam blueprint this chapter aligns to. At
\url{https://apstudents.collegeboard.org} (verified 2026-10-11).
\item D.~Morin, \emph{Introduction to Classical Mechanics}, ch.~5 ---
momentum and rockets at the harder level; ISBN 978-0521876223.
\item The Physics Classroom, \emph{Momentum and Its Conservation} ---
readable remediation for the whole chapter, at
\url{https://www.physicsclassroom.com/class/momentum} (verified 2026-10-11).
\item Khan Academy, \emph{Impacts and linear momentum} --- drills matching
\S\ref{sec:mom-p}--\S\ref{sec:mom-2d}, at
\url{https://www.khanacademy.org/science/physics} (verified 2026-10-11).
\end{itemize}

\section*{Additional exercises}
\addcontentsline{toc}{section}{Additional exercises}

\begin{exercise}[easy]
A constant $500$\,N force acts for $0.2$\,s on a $2$\,kg mass initially at
rest. Find the impulse and final velocity. (AP-C)
\end{exercise}
\begin{exercise}[easy]
A $1200$\,kg car moving at $15$\,m/s is stopped by an identical parked car
they stick to. Find the shared velocity. (F=ma)
\end{exercise}
\begin{exercise}[easy]
Find the CM of three masses: $1$\,kg at $x=0$, $2$\,kg at $x=3$,
$3$\,kg at $x=6$. (AP-C)
\end{exercise}
\begin{exercise}[easy]
Explain why momentum is \emph{not} conserved while a thrown ball is in the
air, in one sentence using the boundary rule. (AP-C)
\end{exercise}
\begin{exercise}[easy]
A cannon fires a $5$\,kg shell at $400$\,m/s; the cannon recoils at
$2$\,m/s. Find the cannon's mass. (F=ma)
\end{exercise}
\begin{exercise}[medium]
Ballistic pendulum, momentum stage only: a $10$\,g bullet at
$v$ stops inside a $2$\,kg block, which then holds together. Find the
block's post-impact speed if $v=300$\,m/s. (The swing is ch4's problem.)
(AP-C)
\end{exercise}
\begin{exercise}[medium]
From the impulse graph of figure~\ref{fig:impulse}: if the same $J$ were
delivered as a \emph{constant} force over $20$\,ms, what would $F$ be? How
does the peak change? (F=ma)
\end{exercise}
\begin{exercise}[medium]
Water jet on a bucket: the hose of example~\ref{ex:jet} hits a bucket
moving away at $5$\,m/s. Redo the thrust using relative flux. (USAPhO-lite)
\end{exercise}
\begin{exercise}[medium]
A $70$\,kg astronaut throws a $2$\,kg wrench at $6$\,m/s while floating.
Find her recoil, then her drift over the $30$\,s EVA repair. (F=ma)
\end{exercise}
\begin{exercise}[medium]
In the CM frame: transform the equal-mass elastic collision of the code
listing into the CM frame, show the velocities merely \emph{reverse sign},
and transform back. (F=ma)
\end{exercise}
\begin{exercise}[medium]
Two pucks collide on air hockey as in example~\ref{ex:glance}, but now the
second puck emerges at $45^{\circ}$ below the original direction with speed
$3$\,m/s. Recompute puck 1's exit. (AP-C)
\end{exercise}
\begin{exercise}[medium]
A rocket has $u=2.5$\,km/s and needs $\Delta v=6$\,km/s in one stage.
Find $m_0/m_f$, then the propellant fraction. (F=ma)
\end{exercise}
\begin{exercise}[hard]
A uniform chain of length $L$ and mass $m$ hangs with its lower end just
touching a table and is released. Show that when a length $x$ has landed,
the table's reaction is $R=\frac{3mx}{L}g$ --- the flux term triples the
weight. (APhO-style)
\end{exercise}
\begin{exercise}[hard]
Two-stage optimisation: with total propellant fraction fixed at $90\%$ and
$u=3$\,km/s, sketch how splitting it $40/50$ vs $60/30$ between stages
changes total $\Delta v$ --- set up, don't solve, the constrained optimum.
(USAPhO-style)
\end{exercise}
\begin{exercise}[medium]
Superball bounce: a ball hits the floor at $v$ and rebounds with
restitution $e=0.9$. Find the impulse delivered and interpret the loss.
(F=ma)
\end{exercise}
